\section{Research Gap and Contribution}

From this review of related work, two major research gaps emerge, which define the contribution of this thesis.

First, while domain-adaptive pre-training is a well-established paradigm for English legal \ac{NLP} \cite{chalkidis-etal-2020-legal, niklaus_2024_multilegalpile}, the German legal domain remains under-explored in this specific context. Foundational work has demonstrated the utility of static word embeddings for query expansion \cite{landthaler_improving_2019} and established general \ac{LIR} benchmarks \cite{bönisch2025finding,wrzalik_gerdalir_2021}. For German tax law specifically, recent advancements have focused almost exclusively on the generative component of legal \ac{AI} systems \cite{wind2026steuerllm, luketina2025using}. Consequently, no prior work has systematically evaluated the suitability of modern, contextualised Transformer-based embedding models for the retrieval stage.

Second, while the intrinsic geometric properties of embedding spaces are well characterised for general-domain models \cite{ethayarajh-2019-contextual,wang-isola-2020-contrastive} and post-hoc corrections have been shown to improve representational quality \cite{Li_2020_Sentence,su2021whitening}, their application to specialised legal domains is entirely absent. The question of whether these geometric metrics carry predictive information about downstream retrieval performance in complex domains and can thus serve as lightweight diagnostic proxies in settings where extrinsic human evaluation is costly, remains open.

This thesis addresses both gaps. It provides an empirical characterisation of how general-purpose and domain-adapted embedding models perform specifically in the German tax law domain, combining intrinsic geometric analysis with task-based extrinsic retrieval evaluation within an operational \ac{RAG} framework. Furthermore, it systematically compares various optimisation strategies, including post-hoc geometric transformations and domain-specific fine-tuning, to investigate the extent to which intrinsic geometric metrics can serve as reliable indicators for retrieval quality. In doing so, it complements the generator-focused work of \citet{wind2026steuerllm} and \citet{luketina2025using} by providing the retriever-focused perspective that has been absent from the German tax law \ac{AI} landscape.


